\documentclass[]{article}

\usepackage{amsmath}
\usepackage{multirow}
\usepackage{natbib}
\usepackage{graphicx} 


\begin{document}

\subsection{Dataset}
The analysis is performed on a suite of ten synthetic galaxy catalogs, generated from a known Halo Occupation Distribution (HOD) model. These catalogs provide the three-dimensional positions, luminosities, and host dark matter halo properties for each galaxy. For our analysis, we focus on the population of satellite galaxies residing in massive halos. The final aggregated dataset used for the likelihood estimation comprises over 68,000 satellite galaxies hosted by approximately 20,000 dark matter halos with virial masses $M_{\text{vir}} \ge 10^{13} M_\odot/h$. For each satellite galaxy, we compute its radial distance, $r$, from the center of its assigned host halo.

\subsection{The Mass-Conditioned Spatial Point Process Model}
We model the three-dimensional spatial distribution of satellite galaxies as a mass-conditioned Neyman-Scott process. In this framework, the centers of dark matter halos act as the ``parent'' points of the process, and the satellite galaxies are the ``offspring'' points distributed around them. The model is defined by two primary components: the mean number of satellites per halo and their spatial distribution.

The mean number of satellite galaxies, $\langle N_{\text{sat}} \rangle$, within a halo of mass $M$ is described by a power-law relation characteristic of HOD models:
\begin{equation}
\langle N_{\text{sat}}(M) \rangle = \left( \frac{M}{M_{\text{sat}}} \right)^{\alpha_{\text{sat}}}
\end{equation}
where $M_{\text{sat}}$ is the characteristic mass scale for a halo to host one satellite, and $\alpha_{\text{sat}}$ is the power-law index.

The spatial positions of these satellites relative to their host halo center are drawn from a radial probability density function, $f(r|M)$, which acts as the spatial kernel of the Neyman-Scott process. We model this kernel using an exponential profile, scaled by the halo's virial radius, $R_{\text{vir}}$:
\begin{equation}
f(r|M) \propto \exp\left(-\alpha_c(M) \frac{r}{R_{\text{vir}}}\right)
\end{equation}
where $\alpha_c(M)$ is a mass-dependent concentration parameter that dictates how centrally concentrated the satellites are. The full intensity function for satellites in a halo of mass $M$ is then given by $\lambda(r | M, \theta) = \langle N_{\text{sat}}(M) \rangle \cdot f(r|M)$, where $\theta$ is the set of model parameters $\{M_{\text{sat}}, \alpha_{\text{sat}}, \alpha_c(M)\}$.

\subsection{Parameter Estimation and Model Selection}
We infer the model parameters $\theta$ via Maximum Likelihood Estimation (MLE). The likelihood function is constructed from the positions of individual satellite galaxies, assuming they represent a Poisson point process conditioned on the observed halo population. To isolate the 1-halo regime and mitigate contamination from neighboring halos (the 2-halo term), we apply a hard radial cut, including only satellites with $r < 5$ Mpc/h in the likelihood calculation.

We test two nested hypotheses for the satellite radial concentration. The first is a baseline model where the concentration is constant across all halo masses, i.e., $\alpha_c(M) = \alpha_0$. The second is a more complex model where the concentration is allowed to vary as a power-law of the host halo mass:
\begin{equation}
\alpha_c(M) = \alpha_0 \left( \frac{M}{M_{\text{pivot}}} \right)^{\beta}
\end{equation}
where we fix the pivot mass to $M_{\text{pivot}} = 10^{13} M_\odot/h$.

To formally compare these two models, we employ the Akaike Information Criterion (AIC), defined as:
\begin{equation}
\text{AIC} = 2k - 2\ln(\mathcal{L})
\end{equation}
where $k$ is the number of free parameters in the model and $\mathcal{L}$ is the maximum value of the likelihood function. The model with the lower AIC value is preferred, with a difference $\Delta \text{AIC} > 10$ considered as decisive evidence in favor of the better-performing model.

\subsection{Marked Correlation Function}
To investigate the spatial segregation of galaxies as a function of their intrinsic properties, we employ a marked correlation function, using galaxy luminosity as the mark. The marked correlation function, $M(r)$, is defined as the mean product of the marks (luminosities) for all pairs of galaxies separated by a distance $r$.

To isolate the correlation between spatial position and luminosity from the underlying spatial clustering of galaxies and the overall luminosity distribution, we normalize $M(r)$ by a shuffled counterpart, $M_{\text{shuffled}}(r)$. This shuffled function is computed by randomly reassigning the luminosity marks among the fixed observed galaxy positions and then recalculating the mean mark product for pairs at separation $r$. The ratio $M(r)/M_{\text{shuffled}}(r)$ provides a direct measure of luminosity segregation. A ratio greater than unity at a given scale indicates that galaxies separated by that distance tend to be more luminous than average, while a ratio less than unity indicates they are less luminous.

\subsection{Residual Analysis}
We evaluate the limitations of our 1-halo Neyman-Scott model by performing a residual analysis in the transition region between 1-halo and 2-halo dominance. We compare the observed number of galaxies, $O$, to the number expected by our best-fit model, $E$, in the radial range of $5 < r < 10$ Mpc/h. This region was excluded from the likelihood fit. The fractional residual, $(O-E)/E$, quantifies the model's breakdown. A large positive residual indicates a significant source of galaxy clustering not captured by the isolated halo model, namely the contribution from neighboring, correlated halos (the 2-halo term).

\end{document}
                